Indeed, take @@M0@@ with @@M1@@, and @@M2@@, so that there is an exact sequence
@@M3@@
by 1.11 and 1.3. Similarly, taking 6.2 into account, take @@M4@@ with @@M5@@, and @@M6@@.

\sgaStableTarget{sga6.e13.thm.6-4}\textbf{Theorem 6.4.} \emph{Under the conditions of 6.1, suppose that @@M7@@ is a quotient of @@M8@@ with @@M9@@. Let @@M10@@ for @@M11@@, where @@M12@@ is the @@M13@@-th @@M14@@-polynomial (1.9); let @@M15@@, and let @@M16@@. Then}:
\begin{enumerate}[label=(\roman*)]
\item \emph{@@M17@@ is @@M18@@-regular.}
\item \emph{@@M19@@}
\item \emph{@@M20@@ is a @@M21@@-sheaf.}
\end{enumerate}

Indeed, take a section @@M22@@ of @@M23@@ such that we obtain a commutative diagram with exact rows and columns
@@M24@@
